\documentclass{beamer}
\usepackage[T1]{fontenc}
\usepackage[utf8]{inputenc}
\usepackage{textcomp}
\usepackage{eurosym}
\DeclareUnicodeCharacter{00B0}{\textdegree}
\DeclareUnicodeCharacter{20AC}{\euro}
\DeclareUnicodeCharacter{2013}{--}
\DeclareUnicodeCharacter{2014}{---}
\DeclareUnicodeCharacter{2018}{`}
\DeclareUnicodeCharacter{2019}{'}
\DeclareUnicodeCharacter{201C}{``}
\DeclareUnicodeCharacter{201D}{''}
\DeclareUnicodeCharacter{2026}{\ldots{}}
\usepackage[spanish,es-noshorthands]{babel}

% Definición del color corporativo aproximado para Pantone Warm Red C
\definecolor{UMWarmRed}{HTML}{F9423A}

% Tema y esquema de colores
\usetheme{Madrid}
\usecolortheme{rose} % Usamos un esquema de color que pueda combinarse bien

% Configuración de colores usando el color corporativo
\setbeamercolor{title}{bg=UMWarmRed, fg=white}
\setbeamercolor{frametitle}{bg=UMWarmRed, fg=white}
\setbeamercolor{structure}{fg=UMWarmRed}
\setbeamercolor{section number projected}{bg=UMWarmRed,fg=white}
\setbeamercolor{itemize item}{fg=UMWarmRed}
\setbeamercolor{itemize subitem}{fg=UMWarmRed}
\setbeamercolor{enumerate item}{fg=UMWarmRed}

\title[DISEÑO DE EXPERIMENTOS]{METODOLOGÍA Y DISEÑO DE EXPERIMENTOS: \\ PRINCIPIOS ÉTICOS}
\author[Alfonso Rosa García]{Alfonso Rosa García }
\institute[G9]{\textbf{Grupo de Universidades G-9}}

\date{9 y 11 de marzo de 2026 \\ 10:00--13:00}

\begin{document}

\begin{frame}
  \titlepage
\end{frame}

\AtBeginSection[] % Comando para hacer algo al inicio de cada sección
{
  \begin{frame}{Índice}
    \tableofcontents[currentsection] % Muestra el índice, enfocando la sección actual
  \end{frame}
}

\section{Consentimiento y privacidad}
\begin{frame}
    \frametitle{Consentimiento y privacidad}
    \centering
    \textbf{\Large Experimentando éticamente respecto a los participantes}
    \vspace{0.5cm}
\end{frame}

\begin{frame}{Introducción a Problemas Éticos en la Investigación}
\textbf{La Importancia de la Ética en la Investigación}
\begin{itemize}
\item Los problemas éticos han marcado la historia de la investigación científica.
\item Casos de falta de consentimiento informado, engaño y daño a los participantes.
\end{itemize}
\textbf{Consecuencias}
\begin{itemize}
    \item Han impulsado el desarrollo de regulaciones éticas en la investigación.
    \item Buscan proteger a los sujetos de investigación y asegurar la integridad de los estudios.
\end{itemize}
\textbf{Tensiones y desafíos}
\begin{itemize}
    \item Encontrar un equilibrio entre la protección de los participantes y el avance del conocimiento científico.
    \item Riesgos potenciales de limitar excesivamente la investigación vs poner en riesgo el bienestar de los sujetos.
    \item Requiere un análisis cuidadoso de riesgos y beneficios en cada proyecto.
\end{itemize}
\end{frame}

\begin{frame}{Estudios históricos notorios}
    \textbf{Casos emblemáticos}
    \begin{itemize}
        \item \textbf{El estudio de Tuskegee sobre la Sífilis:} Negación deliberada de tratamiento eficaz (incluida la penicilina) a participantes afroamericanos.
        \item \textbf{El Experimento de Milgram:} Medición de la disposición a obedecer autoridades a costa de infligir dolor.
        \item \textbf{El Estudio de la Prisión de Stanford:} Simulación de prisión que llevó a comportamientos abusivos y estrés emocional.
    \end{itemize}

    \textbf{Impacto}
    \begin{itemize}
        \item Destacan la explotación, el racismo sistémico y el trauma psicológico en la investigación.
    \end{itemize}
\end{frame}

\begin{frame}{Ejemplos contemporáneos de dilemas éticos}
\textbf{Manipulación de Información en Redes Sociales:} Un estudio de Facebook alteró los feeds de noticias de los usuarios sin su consentimiento para estudiar cambios en las emociones. Este experimento planteó preguntas sobre el consentimiento informado y el derecho a la privacidad en la investigación psicológica en plataformas digitales.

\textbf{CRISPR y Edición genética:} La utilización de CRISPR-Cas9 para editar genéticamente embriones humanos ha levantado preocupaciones éticas profundas sobre la manipulación genética, las implicaciones a largo plazo para la especie humana y los dilemas morales asociados a la "creación" de seres humanos con características deseables.

\textbf{Acceso y uso de datos personales:} La investigación que utiliza grandes bases de datos personales, como historiales médicos o preferencias de consumo, ha puesto en relieve la necesidad de proteger la privacidad de las personas y evitar el uso indebido de la información personal.

\end{frame}


\begin{frame}{Respuesta y regulaciones éticas}
    \textbf{Desarrollo de normativas éticas}
    \begin{itemize}
        \item \textbf{Declaración de Helsinki:} Principios éticos para la investigación médica.
        \item \textbf{Informe Belmont:} Establece respeto por las personas, beneficencia y justicia.
    \end{itemize}

    \textbf{Impacto en la investigación contemporánea}
    \begin{itemize}
        \item Obligación de obtener consentimiento informado.
        \item Supervisión por Comités de Ética de Investigación (IRBs).
        \item Promoción de transparencia y responsabilidad en la publicación de estudios.
    \end{itemize}
\end{frame}


% Diapositiva 1: Introducción y Principios Generales
\begin{frame}{Declaración de Helsinki: Introducción y Principios Generales}
\begin{itemize}
    \item \textbf{Origen:} Adoptada en 1964 por la Asociación Médica Mundial.
    \item \textbf{Propósito:} Proteger la integridad, dignidad y derechos de los participantes en investigaciones médicas.
    \item \textbf{Fundamentos éticos:} 
    \begin{itemize}
         \item Respeto a la autonomía y autodeterminación.
         \item Principio de beneficencia y no maleficencia.
         \item Prioridad del bienestar del sujeto sobre los intereses de la ciencia y la sociedad.
    \end{itemize}
    \item \textbf{Evolución:} Ha sido revisada en diversas ocasiones para adaptarse a los avances científicos y éticos.
    \item \textbf{Versión actual:} \href{https://www.wma.net/es/policies-post/declaracion-de-helsinki-de-la-amm-principios-eticos-para-las-investigaciones-medicas-en-seres-humanos/}{Declaración de Helsinki de la AMM (2024)}

\end{itemize}
\end{frame}

% Diapositiva 2: Consentimiento Informado y Protección de Participantes
\begin{frame}{Declaración de Helsinki: Consentimiento Informado y Protección de Participantes}
\begin{itemize}
    \item \textbf{Consentimiento Informado:}
    \begin{itemize}
         \item Información clara sobre objetivos, métodos, riesgos y beneficios del estudio.
         \item Obtención del consentimiento de forma libre y voluntaria, permitiendo la revocación en cualquier momento.
    \end{itemize}
    \item \textbf{Protección del Participante:}
    \begin{itemize}
         \item Salvaguarda de la privacidad y confidencialidad de los datos.
         \item Garantizar que la participación no comprometa la integridad física ni mental del sujeto.
    \end{itemize}
\end{itemize}
\end{frame}

% Diapositiva 3: Revisión Ética y Protección de Grupos Vulnerables
\begin{frame}{Declaración de Helsinki: Revisión Ética y Protección de Grupos Vulnerables}
\begin{itemize}
    \item \textbf{Revisión Ética:}
    \begin{itemize}
         \item Todos los estudios deben someterse a la evaluación de un comité de ética independiente.
         \item Supervisión continua del desarrollo del estudio para garantizar el cumplimiento ético.
    \end{itemize}
    \item \textbf{Protección de Grupos Vulnerables:}
    \begin{itemize}
         \item Implementar medidas adicionales para proteger a niños, ancianos, personas con discapacidad, entre otros.
         \item Asegurar que la investigación en estos grupos solo se realice si es imprescindible y con salvaguardias específicas.
    \end{itemize}
\end{itemize}
\end{frame}

% Diapositiva 4: Transparencia, Publicación y Responsabilidad del Investigador
\begin{frame}{Declaración de Helsinki: Transparencia, Publicación y Responsabilidad del Investigador}
\begin{itemize}
    \item \textbf{Transparencia:}
    \begin{itemize}
         \item Registro público de los ensayos clínicos y divulgación de información relevante.
         \item Revelación de conflictos de interés y fuentes de financiamiento.
    \end{itemize}
    \item \textbf{Publicación de Resultados:}
    \begin{itemize}
         \item Obligación de publicar resultados positivos, negativos e inconclusos.
         \item Garantizar la integridad y veracidad de la información difundida.
    \end{itemize}
    \item \textbf{Responsabilidad del Investigador:}
    \begin{itemize}
         \item Compromiso ético en el seguimiento y reporte de los resultados.
         \item Responsabilidad de proveer tratamiento y compensación en caso de daños derivados de la investigación.
    \end{itemize}
\end{itemize}
\end{frame}

\section{Ética en la identificación y difusión de resultados}

\begin{frame}
    \frametitle{Ética en la identificación y difusión de resultados}
    \centering
    \textbf{\Large Reproducibilidad y comportamientos éticos}
    \vspace{0.5cm}
\end{frame}

\begin{frame}{La Crisis de reproducibilidad en la ciencia}
\textbf{¿Qué es la Crisis de reproducibilidad?}
\begin{itemize}
    \item Desde Ioannidis (2005), "Why Most Published Research Findings Are False", se ha encontrado que numerosos resultados de estudios científicos no son reproducibles.
    \item Prinz et al (2011) mostraron que apenas 25\% de los estudios preclínicos publicados pudieron ser validados para continuar.
    \item Begley et al (2012) indicaron que solo pudieron confirmar el 6\% de esos estudios.   
    \item Open Science Collaboration (2015) replicó 100 artículos de 2008 de revistas de psicología de alto impacto y encontraron que solo el 47\% de los resultados originales estaban en el intervalo de confianza del 95\% del tamaño de efecto replicado.
    \item Camerer et al (2016) encontraron que el 39\% de los resultados experimentales publicados en dos de las revistas top de economía no se replicaban.
\end{itemize}
\end{frame}

\begin{frame}{Tarea: Reflexión sobre las causas de la crisis de reproducibilidad}
\textbf{Objetivo:} Identificar y analizar los factores que contribuyen a la crisis de reproducibilidad en la investigación científica.\\[10pt]
\textbf{Instrucciones:}
\begin{enumerate}
    \item Formar grupos de 3-4 personas.
    \item Reflexionar sobre los siguientes aspectos:
    \begin{itemize}
        \item ¿Qué factores del proceso de investigación podrían llevar a resultados irreproducibles?
        \item ¿Cómo influyen la presión por publicar y las prácticas de análisis de datos (como el p-hacking) en la reproducibilidad?
        \item ¿Qué papel juegan las revistas científicas y los revisores en este problema?
    \end{itemize}
    \item Discutir posibles soluciones que podrían implementarse en la comunidad científica para mejorar la reproducibilidad.
\end{enumerate}
\end{frame}

\begin{frame}{La Crisis de reproducibilidad en la ciencia}
\textbf{Causas Principales}
\begin{itemize}
    \item Prácticas de investigación cuestionables, como el p-hacking y la selección selectiva de resultados.
    \item Presiones sobre los investigadores para publicar resultados positivos.
    \item Falta de transparencia y disponibilidad de datos brutos para análisis independientes.
\end{itemize}

\textbf{Impacto en la Ciencia}
\begin{itemize}
    \item Erosión de la confianza pública en la ciencia.
    \item Llamados a la reforma de prácticas de investigación y publicación.
\end{itemize}

\textbf{Respuestas a la Crisis}
\begin{itemize}
    \item Fomento de la ciencia abierta y el preregistro de estudios.
    \item Mayor énfasis en la replicación y validación de resultados por parte de la comunidad científica.
\end{itemize}
\end{frame}


\begin{frame}{Replicación y reproducibilidad en la Investigación experimental}
\textbf{Pilares de la validez científica}
\begin{itemize}
    \item \textbf{Replicación:} Capacidad de obtener resultados al repetir un experimento con mismas condiciones y métodos.
    \item \textbf{Reproducibilidad:} Capacidad de obtener hallazgos consistentes por diferentes investigadores al aplicar los mismos procedimientos.
\end{itemize}

\textbf{Importancia}
\begin{itemize}
    \item Verifica la fiabilidad y solidez de los resultados experimentales.
    \item Esencial para la acumulación de conocimiento científico.
    \item Identifica y corrige posibles errores o sesgos en los estudios originales.
    \item Fomenta la transparencia y la integridad en la investigación.
\end{itemize}

\textbf{Estrategias para promover la replicación y reproducibilidad}
\begin{itemize}
    \item Descripción detallada de los métodos y materiales.
    \item Compartir datos y código de análisis de manera abierta.
    \item Preregistro para distinguir análisis confirmatorios y exploratorios.
\end{itemize}
\end{frame}

\begin{frame}{Multiple testing y control del FWER}
    \textbf{El Desafío de las Pruebas Múltiples}
    \begin{itemize}
        \item Riesgo de falso positivo aumenta con cada prueba estadística adicional.
        \item Importancia de controlar la tasa de error familiar (FWER).
    \end{itemize}

    \textbf{Métodos de Control}
    \begin{itemize}
        \item Procedimiento de Bonferroni y método de Holm para ajustar valores p.
    \end{itemize}

    \textbf{Impacto en la Investigación}
    \begin{itemize}
        \item Asegura que resultados significativos sean genuinos.
        \item Fortalece la confianza en los hallazgos reportados.
    \end{itemize}
\end{frame}

\begin{frame}{Cherry Picking y la Ciencia Abierta}
    \textbf{Contra la selección sesgada de datos}
    \begin{itemize}
        \item Práctica de reportar solo resultados que apoyan hipótesis predeterminadas.
        \item Compromete la reproducibilidad y validez de los estudios.
    \end{itemize}

    \textbf{Mitigando el cherry picking}
    \begin{itemize}
        \item Adopción de prácticas de ciencia abierta y transparencia investigativa.
        \item Publicación de todos los resultados, apoyen o contradigan las hipótesis.
    \end{itemize}

    \textbf{Construyendo una base de conocimiento Sólida}
    \begin{itemize}
        \item Interrelación de pre-registro, control de multiple testing y ciencia abierta.
        \item Promueve una interpretación objetiva de los datos y resultados genuinos.
    \end{itemize}
\end{frame}

\begin{frame}{Pre-registro de estudios científicos}
    \textbf{Aumentando la transparencia}
    \begin{itemize}
        \item Registro público del plan de investigación antes de la recolección de datos.
        \item Combate el p-hacking y el HARKing, fortaleciendo la validez de los resultados.
    \end{itemize}

    \textbf{Beneficios clave}
    \begin{itemize}
        \item Diferencia clara entre hipótesis exploratorias y confirmatorias.
        \item Evita la sobreinterpretación de correlaciones fortuitas.
    \end{itemize}

    \textbf{Herramientas para el pre-registro}
    \begin{itemize}
        \item Plataformas como el Registro de Ensayos Clínicos de los NIH, OSF o AsPredicted.
        \item \textbf{Ejemplos:} \url{https://osf.io/93aus/} o \url{https://aspredicted.org/2zzm-5r7q.pdf}
    \end{itemize}
\end{frame}

\end{document}








